\section{Latihan}
\label{sec:latihan-bab17}

Latihan pada bab ini bersifat reflektif dan integratif. Tidak ada
jawaban tunggal yang benar; yang dinilai adalah kedalaman dan
kejujuran analisis Anda, serta ketepatan bukti yang ditunjuk.

\begin{enumerate}
    \item \textbf{Integrasi konsep.} Jelaskan bagaimana
          Bab~\ref{chap:segmentasi}, Bab~\ref{chap:mode-pengalamatan},
          dan Bab~\ref{chap:prosedur} saling terkait ketika sebuah
          prosedur menerima parameter melalui stack. Urutkan
          konsepnya dan tunjukkan di mana setiap bab berperan.

    \item \textbf{Analisis kode lintas bab.} Diberikan sebuah
          program yang membaca nama pengguna, menghitung panjangnya,
          lalu mencetak salam. Identifikasi bab mana yang menyediakan
          konsep untuk setiap bagian program tersebut.

    \item \textbf{Penilaian diri.} Isi lembar penilaian diri pada
          Tabel~\ref{tab:self-assessment}. Untuk setiap CPMK, sebutkan
          satu bukti konkret berupa nama berkas atau nomor tugas.
          Klaim tanpa bukti tidak dihitung.

    \item \textbf{Analisis kesalahan.} Pilih satu bug yang pernah
          Anda alami sepanjang semester. Jelaskan: gejala yang
          muncul, hipotesis pertama Anda, langkah yang membuktikan
          atau menggugurkan hipotesis itu, dan penyebab sebenarnya.

    \item \textbf{Perbandingan pendekatan.} Ambil satu program yang
          Anda tulis pada Bab~\ref{chap:kontrol}. Tuliskan ulang
          bagian perulangannya dengan pendekatan berbeda (misalnya
          mengganti \asm{LOOP} dengan \asm{DEC} dan \asm{JNZ}), lalu
          bandingkan jumlah instruksi dan keterbacaannya.

    \item \textbf{Evaluasi rubrik.} Baca kembali rubrik proyek pada
          Tabel~\ref{tab:rubrik-proyek}. Untuk setiap kriteria,
          tuliskan satu hal konkret yang akan Anda lakukan pada
          proyek akhir agar memperoleh skor 4.

    \item \textbf{Rencana tindak lanjut.} Berdasarkan
          Tabel~\ref{tab:tindak-lanjut}, pilih satu jalur lanjutan,
          lalu susun rencana belajar tiga bulan pertama: topik,
          sumber, dan bukti yang akan Anda hasilkan.

    \item \textbf{Kritik buku.} Sebutkan satu bagian buku ini yang
          menurut Anda paling kurang jelas, dan usulkan cara
          perbaikannya. Sertakan nomor bab dan bagian.
\end{enumerate}

\begin{refleksibox}
Jurnal refleksi akhir semester. Jawab dengan jujur dan sertakan
bukti.
\begin{enumerate}
    \item Kompetensi yang paling berkembang pada diri saya adalah
          \dots, dengan bukti \dots
    \item Kompetensi yang masih paling lemah adalah \dots, dan
          rencana perbaikan saya adalah \dots
    \item Kebiasaan belajar yang paling membantu saya adalah \dots
    \item Kebiasaan belajar yang paling merugikan saya adalah \dots
    \item Setelah mempelajari level instruksi, hal yang paling
          mengubah cara saya memandang program adalah \dots
\end{enumerate}
\end{refleksibox}
